\documentclass[a4paper,12pt]{article}
\usepackage[utf8]{inputenc}
\usepackage[margin=1in]{geometry}
\usepackage{xcolor}
\usepackage{hyperref}
\usepackage{titlesec}
\usepackage{enumitem}
\usepackage{fontenc}
\usepackage{microtype}
\usepackage[normalem]{ulem}

% Couleurs personnalisées
\definecolor{maincolor}{HTML}{C60000}
\definecolor{graytext}{HTML}{555555}

% Configuration des liens
\hypersetup{
	colorlinks=true,
	linkcolor=maincolor,
	urlcolor=maincolor
}

% Formatage des titres
\titleformat{\section}%
{\normalfont\large\bfseries\color{maincolor}}{}{0em}{}[\color{maincolor}\titlerule]

% Formatage personnalisé des éléments de liste
\setlist[itemize]{topsep=0pt, left=1em, itemsep=3pt}

\begin{document}
	
	% En-tête
	\begin{center}
		\textbf{\Huge MOHAMED EL BACHIR}\\
		\vspace{2mm}
		\textbf{\Large BOUBA NGANADAKOUA}
	\end{center}
	
	\noindent
	\textbf{Date et lieu de naissance :} 02 août 2002, Garoua\\
	\textbf{Nationalité :} Camerounaise\\
	\textbf{Téléphone :} (+237) 698 34 06 64\\
	\textbf{Email :} \href{mailto:mohamedelbachirboubanganadakou@gmail.com}{mohamedelbachirboubanganadakou@gmail.com}\\
	\textbf{GitHub :} \href{https://github.com/mohamedelbachir}{github.com/mohamedelbachir}\\
	\textbf{Portfolio :} \href{https://bachdev.vercel.app}{bachdev.vercel.app}\\
	\textbf{LinkedIn :} \href{https://www.linkedin.com/in/mohamed-el-bachir/}{linkedin.com/in/mohamed-el-bachir/}
	
	\vspace{5mm}
	
	% Présentation
	\section*{Présentation}
	Je suis titulaire d'un Master 2 en Informatique, spécialisé dans l'architecture des systèmes logiciels et les environnements distribués. Passionné par la programmation et l'innovation technologique, je souhaite mettre mes compétences au service de projets ayant un impact positif, en apportant des solutions modernes et performantes.
	
	\vspace{5mm}
	
	% Formation
	\section*{Formation}
	\begin{itemize}
		\item \textbf{Master 2 en Informatique} – Spécialité Système Logiciel et Environnements Distribués (\textit{SLED}) : Thème : \textit{Deep Analysis and Prediction of Chikungunya Using Ensemble Regression Approach}, 2024, mention \textit{Très Bien}.
		\item \textbf{Master 1 Recherche en Informatique} – Spécialité \textbf{SLED}, Université de Ngaoundéré, 2023, mention \textit{Assez Bien}, Major de promotion.
		\item \textbf{Licence en Informatique}, Université de Ngaoundéré, 2022, mention \textit{Assez Bien}.
		\item \textbf{Baccalauréat Série C}, 2019, mention \textit{Assez Bien}.
	\end{itemize}
	
	\vspace{5mm}
	
	% Langues
	\section*{Langues}
	\begin{itemize}
		\item \textbf{Français} : Courant
		\item \textbf{Anglais} : Niveau B1
	\end{itemize}
	
	% Compétences
	\section*{Compétences Techniques}
	\begin{itemize}
		\item \textbf{Développement Frontend} : ReactJS, Styled Components, Sass, TailwindCSS, NextJS, AstroJS, React Native.
		\item \textbf{Développement Backend} : NodeJS, ExpressJS, Firebase, Drizzle ORM, PostgreSQL.
		\item \textbf{Bibliothèques et Frameworks} : SDL, ImGUI, Mantine UI, Shadcn UI, Material UI.
		\item \textbf{Langages de Programmation} : C, C++, JavaScript, TypeScript, Python.
		\item \textbf{Outils} : Git, VSCode, CMake.
		\item \textbf{Déploiement} : Netlify, Vercel, WebHostApp, Render.
	\end{itemize}
	
	% Expérience Professionnelle
	\section*{Expérience Professionnelle}
	\begin{itemize}
		\item \textbf{2024 – Présent} : Stage Professionnel en tant que \textbf{Développeur Frontend} chez \textbf{Sygalin Sas CANAL +}, Ngaoundéré. Développement d'une application de gestion d'association (AssoEase) en utilisant \textbf{React}, \textbf{TypeScript}, \textbf{Mantine UI}.
		\item \textbf{2019} : Attestation de fin de formation en \textbf{Initiation à la Maintenance et aux Caméras de Surveillance} à la \textbf{GI} (Généralité Informatique), Garoua, mention \textit{Bien}.
	\end{itemize}
	
	% Certifications
	\section*{Certifications}
	\begin{itemize}
		\item \textbf{C++} – Sololearn : \href{https://www.sololearn.com/en/certificates/CT-X8YY1SDM}{Référence}.
		\item \textbf{Building Blog using MERN Stack} – Udemy : \href{https://www.udemy.com/certificate/UC-e01be5eb-a327-4f45-b5f2-73cc5fcf5d0c/}{Référence}.
		\item \textbf{Latex} – AUF : \href{https://www.linkedin.com/in/mohamed-el-bachir/}{Référence}.
		\item \textbf{CITS} – Cameroon International Tech Summit : Participant au hackathon national en tant que dévéloppeur frontend et \textbf{Chef d'équipe} d'une équipe de \textbf{5 personnes}, \href{https://www.linkedin.com/in/mohamed-el-bachir/}{Référence}.
	\end{itemize}
	
	% Centres d'Intérêt
	\section*{Centres d'Intérêt}
	Programmation, Sciences, Animaux, Anime, Jeux Vidéo, Bandes Dessinées.
	
	% Projets
	\section*{Projets}
	\begin{itemize}
		\item \textbf{Scheduler} : Logiciel de visualisation de diagrammes de GANT pour l’ordonnancement des processus dans les systèmes UNIX. Développé en \textbf{C++}.
		\begin{itemize}
			\item Gestion de thèmes.
			\item Implémentation des algorithmes d'ordonnancement : Round Robin, FCFS, SJF, SRTF.
			\item Documentation complète.
		\end{itemize}
		\textbf{Outils} : C++, SDL, ImGui, SDL\_TTF, SDL\_Image, TypeScript, Astro.\\
		\href{https://github.com/mohamedelbachir/Scheduler}{Voir le projet.}
		
		\item \textbf{U-sentry} : Plateforme de gestion d'alertes dans le milieu universitaire. Développée avec \textbf{React}, \textbf{Expo}, \textbf{TypeScript}, \textbf{Redux}, \textbf{Supabase}, \textbf{React Native}.
		\begin{itemize}
			\item Gestion des posts et des notifications.
			\item Interface utilisateur mobile avec parcours personnalisé et favoris.
			\item Mises à jour en temps réel avec Supabase.
		\end{itemize}
		\href{https://bachdev.vercel.app/projects/u-sentry}{Voir le projet.}
		
		\item \textbf{Chair Rental} : Plateforme de réservation de chaises pour événements au Cameroun (en cours).
		\begin{itemize}
			\item Gestion du panier et suivi des produits.
			\item Tableau de bord administrateur.
		\end{itemize}
		\textbf{Outils} : React, TypeScript, Redux.\\
		\href{https://chair-rental.vercel.app/}{Voir le projet.}
		
		\item \textbf{Généralité Informatique} : E-commerce de produits technologiques, y compris caméras et fournitures. 
		\begin{itemize}
			\item Affichage des produits et gestion du catalogue.
			\item Interface de type landing page et tableau de bord administrateur.
		\end{itemize}
		\textbf{Outils} : React, JavaScript, Firebase, Material UI, Sass.\\
		\href{https://generaliteinformatique.vercel.app}{Voir le projet.}
		
		\item \textbf{Woconnex} : Plateforme développée lors du hackathon CITS 2024 sur la disparité de genre dans les secteurs technologiques (en cours).
		\begin{itemize}
			\item Gestion des profils, du feed, et de l’abonnement.
		\end{itemize}
		\textbf{Outils} : React, TypeScript, Redux.\\
		\href{https://woconnex.vercel.app/dashboard}{Voir le projet.}
		
		\item \textbf{Jumper} : Jeu éducatif en \textbf{C++} pour apprendre les bases de la programmation.
		\begin{itemize}
			\item Gestion de l’écran, des sons et des contrôles.
		\end{itemize}
		\textbf{Outils} : C++, SDL.\\
		\href{https://github.com/mohamedelbachir/JumperGame}{Voir le projet.}
	\end{itemize}
	
	\vspace{5mm}
	Pour découvrir plus de projets, visitez : \href{https://bachdev.vercel.app/projects}{https://bachdev.vercel.app/projects}
	
\end{document}
